\documentclass[10pt,oneside,openany]{memoir}
%%%%%%%%%%%%%%%%%%%%%%%%%%%%%%%%%%%%%%%%%%%%%%%%%%%%%%%%%%%%%%%%%%%%%%%%%%%%%%%%%%%%%
% Font options
%%%%%%%%%%%%%%%%%%%%%%%%%%%%%%%%%%%%%%%%%%%%%%%%%%%%%%%%%%%%%%%%%%%%%%%%%%%%%%%%%%%%%

\iffalse
%font style 1
\usepackage[p,osf]{scholax}
\usepackage{amsmath,amsthm,amssymb}
\usepackage[scaled=1.075,ncf,vvarbb]{newtxmath}
\fi

%font style2
\usepackage[T1]{fontenc}
\usepackage{amsmath}
\usepackage[fixamsmath]{mathtools}  % Extension to amsmath
\usepackage{amsthm}
\usepackage{amssymb}
\renewcommand*{\mathbf}[1]{\mathbb{#1}}
\renewcommand{\qedsymbol}{$\blacksquare$}

%font style 3
%\usepackage{newpxtext}
%\usepackage{eulerpx}
%\usepackage{mathrsfs}
%\usepackage{eucal}

\usepackage[uprightgreeks,light]{kpfonts}
\usepackage{datetime}
%%%%%%%%%%%%%%%%%%%%%%%%%%%%%%%%%%%%%%%%%%%%%%%%%%%%%%%%%%%%%%%%%%%%%%%%%%%%%%%%%%%%%
% Page layout and headers
%%%%%%%%%%%%%%%%%%%%%%%%%%%%%%%%%%%%%%%%%%%%%%%%%%%%%%%%%%%%%%%%%%%%%%%%%%%%%%%%%%%%%


\usepackage[margin=1in]{geometry}
\usepackage{graphicx} % needed for \scalebox
\usepackage{fancyhdr}
\usepackage{datetime}

\pagestyle{fancy}
\fancyhf{}
\cfoot{\scriptsize\thepage}
\fancyhead[R]{\scalebox{0.8}{\nouppercase{\rightmark}}}
\fancyhead[L]{\scalebox{0.8}{\nouppercase{\leftmark}}}

\fancypagestyle{plain}{%
    \fancyhf{}
    \cfoot{\scriptsize\thepage}
    \fancyhead[R]{\scalebox{0.8}{\phantom{a}}}
    \fancyhead[L]{\scalebox{0.8}{Applied Analysis}}
}

\renewcommand{\headrulewidth}{0pt}
\renewcommand{\footrulewidth}{0pt}

%%%%%%%%%%%%%%%%%%%%%%%%%%%%%%%%%%%%%%%%%%%%%%%%%%%%%%%%%%%%%%%%%%%%%%%%%%%%%%%%%%%%%
% Theorem environments
%%%%%%%%%%%%%%%%%%%%%%%%%%%%%%%%%%%%%%%%%%%%%%%%%%%%%%%%%%%%%%%%%%%%%%%%%%%%%%%%%%%%%

\usepackage{keytheorems}

\newkeytheoremstyle{def}{
    inherit-style=definition,
    spaceabove=1.5em,
    spacebelow=1.5em,
    postheadspace=0.5em,
}

\newkeytheoremstyle{thm}{
    spaceabove=1.5em,
    spacebelow=1.5em,
    postheadspace=0.5em,
}

\newkeytheorem{theorem}[
    name = Theorem,
    style=thm,
    numberwithin = section,
]

\newkeytheorem{theorem*}[
    name=Theorem,
    style=thm,
    numbered=false
]

\newkeytheorem{proposition}[
    name = Proposition,
    style=thm,
    numberwithin = section,
    sibling=theorem
]

\newkeytheorem{lemma}[
    name = Lemma,
    style=thm,
    numberwithin = section,
    sibling=theorem
]

\newkeytheorem{corollary}[
    name = Corollary,
    style=thm,
    numberwithin = section,
    sibling=theorem
]

\newkeytheorem{definition}[
    name = Definition,
    style=def,
    numberwithin = section,
    style = definition
]

\newkeytheorem{example}[
    name = Example,
    style=def,
    numberwithin = section,
    style = definition
]

\newkeytheorem{remark}[
    name=Remark,
    style=remark,
    numbered=false
]

\newkeytheorem{question}[
    name=Question,
    style=remark,
    numbered=false
]

\newkeytheorem{exercise}[
    name = Exercise,
    style=thm
]
\newkeytheorem{problem}[
    name = Problem,
    style=thm
]

\newkeytheorem{axiom}[name = Axiom]

%%%%%%%%%%%%%%%%%%%%%%%%%%%%%%%%%%%%%%%%%%%%%%%%%%%%%%%%%%%%%%%%%%%%%%%%%%%%%%%%%%%%%
% Section and chapter formatting
%%%%%%%%%%%%%%%%%%%%%%%%%%%%%%%%%%%%%%%%%%%%%%%%%%%%%%%%%%%%%%%%%%%%%%%%%%%%%%%%%%%%%

\usepackage{enumitem}
\usepackage[explicit]{titlesec}

\titleformat{\chapter}[display]
    {\bfseries\LARGE\raggedright}
    {Chapter {\thechapter}}
    {.1ex minus .1ex}
    {\Large #1}

\titlespacing{\chapter}
    {3pc}{*3}{40pt}[3pc]

\titleformat{\section}[block]
    {\normalfont\bfseries\large}
    {}{0pt}
    {\fbox{\strut \S\ \thesection.\ #1}}

\titlespacing{\section}
    {0pt}{3ex plus .1ex minus .2ex}{3ex plus .1ex minus .2ex}

\setlength{\fboxsep}{4pt}    % padding inside the box
\setlength{\fboxrule}{0.6pt} % border thickness

\titleformat{name=\section,numberless}[block]
    {\normalfont\bfseries\large\raggedright}
    {}{0pt}
    {\fbox{\strut #1}}

%%%%%%%%%%%%%%%%%%%%%%%%%%%%%%%%%%%%%%%%%%%%%%%%%%%%%%%%%%%%%%%%%%%%%%%%%%%%%%%%%%%%%
% Graphics, symbols, and miscellaneous packages
%%%%%%%%%%%%%%%%%%%%%%%%%%%%%%%%%%%%%%%%%%%%%%%%%%%%%%%%%%%%%%%%%%%%%%%%%%%%%%%%%%%%%

\usepackage{tikz}
\usepackage{tikz-cd}
\usepackage{float}
\usepackage{wasysym}
\usepackage{hyperref}
\usepackage{ulem}
\usepackage{pgfplots}
\usetikzlibrary{patterns}
\usetikzlibrary{positioning,arrows.meta}
\usetikzlibrary{decorations.pathreplacing}
\usetikzlibrary{patterns}
\usepgfplotslibrary{fillbetween}
\pgfplotsset{compat=1.18}
\usepackage{etoolbox}
\usepackage{bm}

%%%%%%%%%%%%%%%%%%%%%%%%%%%%%%%%%%%%%%%%%%%%%%%%%%%%%%%%%%%%%%%%%%%%%%%%%%%%%%%%%%%%%
% Final Formatting/Extra commands
%%%%%%%%%%%%%%%%%%%%%%%%%%%%%%%%%%%%%%%%%%%%%%%%%%%%%%%%%%%%%%%%%%%%%%%%%%%%%%%%%%%%%

\newcommand{\homeworktitle}[3]{%
    \begin{center}
        {\large #1 \\[0.1in] #2}\\[.175in]
        {Name:} {\underline{#3\hspace*{2in}}}\\[0.15in]
    \end{center}
    \vspace{4pt}
}
\newdateformat{specialdate}{\THEYEAR\ \monthname\ \THEDAY}

%%%%%%%%%%%%%%%%%%%%%%%%%%%%%%%%%%%%%%%%%%%%%%%%%%%%%%%%%%%%%
%%%%%%%%%%%%%%%%%%%%%%%%%%%%%%%%%%%%%%%%%%%%%%%%%%%%%%%%%%%%%

% ============================================================
% makros.tex — reorganized, full restored version
% ============================================================

% ============================================================
% Sha / Cyrillic support notes
% ============================================================

%to make the correct symbol for Sha
%\newcommand\cyr{%
%\renewcommand\rmdefault{wncyr}%
%\renewcommand\sfdefault{wncyss}%
%\renewcommand\encodingdefault{OT2}%
%\normalfont \selectfont} \DeclareTextFontCommand{\textcyr}{\cyr}

% ============================================================
% Math operators
% ============================================================

\DeclareMathOperator{\ab}{ab}
\DeclareMathOperator{\ad}{ad}
\DeclareMathOperator{\adj}{adj}
\DeclareMathOperator{\alg}{alg}
\DeclareMathOperator{\Alt}{Alt}
\DeclareMathOperator{\Ann}{Ann}
\DeclareMathOperator{\arith}{arith}
\DeclareMathOperator{\Aut}{Aut}
\DeclareMathOperator{\Be}{B}
\DeclareMathOperator{\Bd}{Bd}
\DeclareMathOperator{\card}{card}
\DeclareMathOperator{\Char}{char}
\DeclareMathOperator{\csp}{csp}
\DeclareMathOperator{\codim}{codim}
\DeclareMathOperator{\coker}{coker}
\DeclareMathOperator{\coh}{H}
\DeclareMathOperator{\compl}{compl}
\DeclareMathOperator{\conj}{conj}
\DeclareMathOperator{\cont}{cont}
\DeclareMathOperator{\Cov}{Cov}
\DeclareMathOperator{\crys}{crys}
\DeclareMathOperator{\Crys}{Crys}
\DeclareMathOperator{\cusp}{cusp}
\DeclareMathOperator{\diag}{diag}
\DeclareMathOperator{\diam}{diam}
\DeclareMathOperator{\Dom}{Dom}
\DeclareMathOperator{\disc}{disc}
\DeclareMathOperator{\dist}{dist}
\DeclareMathOperator{\Div}{div}
\DeclareMathOperator{\dR}{dR}
\DeclareMathOperator{\Eis}{Eis}
\DeclareMathOperator{\End}{End}
\DeclareMathOperator{\ev}{ev}
\DeclareMathOperator{\eval}{eval}
\DeclareMathOperator{\Eq}{Eq}
\DeclareMathOperator{\Ext}{Ext}
\DeclareMathOperator{\Fil}{Fil}
\DeclareMathOperator{\Fitt}{Fitt}
\DeclareMathOperator{\Frob}{Frob}
\DeclareMathOperator{\G}{G}
\DeclareMathOperator{\Gal}{Gal}
\DeclareMathOperator{\GL}{GL}
\DeclareMathOperator{\Gr}{Gr}
\DeclareMathOperator{\Graph}{Graph}
\DeclareMathOperator{\GSp}{GSp}
\DeclareMathOperator{\GUn}{GU}
\DeclareMathOperator{\Hom}{Hom}
\DeclareMathOperator{\id}{id}
\DeclareMathOperator{\Id}{Id}
\DeclareMathOperator{\Ik}{Ik}
\DeclareMathOperator{\IM}{Im}
\DeclareMathOperator{\Image}{im}
\DeclareMathOperator{\Ind}{Ind}
\DeclareMathOperator{\Inf}{inf}
\DeclareMathOperator{\ins}{ins}
\DeclareMathOperator{\inv}{inv}
\DeclareMathOperator{\Isom}{Isom}
\DeclareMathOperator{\J}{J}
\DeclareMathOperator{\Jac}{Jac}
\DeclareMathOperator{\lcm}{lcm}
\DeclareMathOperator{\length}{length}
\DeclareMathOperator*{\limit}{limit}
\DeclareMathOperator{\Log}{Log}
\DeclareMathOperator{\M}{M}
\DeclareMathOperator{\Mat}{Mat}
\DeclareMathOperator{\N}{N}
\DeclareMathOperator{\Nm}{Nm}
\DeclareMathOperator{\NIk}{N-Ik}
\DeclareMathOperator{\NSK}{N-SK}
\DeclareMathOperator{\new}{new}
\DeclareMathOperator{\obj}{obj}
\DeclareMathOperator{\old}{old}
\DeclareMathOperator{\ord}{ord}
\DeclareMathOperator{\Or}{O}
\DeclareMathOperator{\op}{op}
\DeclareMathOperator{\out}{out}
\DeclareMathOperator{\PGL}{PGL}
\DeclareMathOperator{\PGSp}{PGSp}
\DeclareMathOperator{\rank}{rank}
\DeclareMathOperator{\Ran}{Ran}
\DeclareMathOperator{\rel}{Rel}
\DeclareMathOperator{\Real}{Re}
\DeclareMathOperator{\RES}{res}
\DeclareMathOperator{\Res}{Res}
%\DeclareMathOperator{\Sha}{\textcyr{Sh}}
\DeclareMathOperator{\Sel}{Sel}
\DeclareMathOperator{\semi}{ss}
\DeclareMathOperator{\sgn}{sign}
\DeclareMathOperator{\SK}{SK}
\DeclareMathOperator{\SL}{SL}
\DeclareMathOperator{\SO}{SO}
\DeclareMathOperator{\Sp}{Sp}
\DeclareMathOperator{\Span}{span}
\DeclareMathOperator{\Spec}{Spec}
\DeclareMathOperator{\spin}{spin}
\DeclareMathOperator{\st}{st}
\DeclareMathOperator{\St}{St}
\DeclareMathOperator{\SUn}{SU}
\DeclareMathOperator{\supp}{supp}
\DeclareMathOperator{\Sup}{sup}
\DeclareMathOperator{\Sym}{Sym}
\DeclareMathOperator{\Tam}{Tam}
\DeclareMathOperator{\tors}{tors}
\DeclareMathOperator{\tr}{tr}
\DeclareMathOperator{\Tr}{Tr}
\DeclareMathOperator{\un}{un}
\DeclareMathOperator{\Un}{U}
\DeclareMathOperator{\val}{val}
\DeclareMathOperator{\vol}{vol}

% ============================================================
% General aliases and short macros
% ============================================================

\newcommand{\absgal}{\G_{\bbQ}}
\newcommand{\Loc}{L_{\operatorname{loc}}}
\renewcommand{\k}{\kappa}
\newcommand{\Ff}{F_{f}}
%\newcommand{\ts}{\,^{t}\!}
\newcommand{\lat}{\mathscr{L}}
\newcommand{\ov}[1]{\overline{#1}}
\newcommand{\up}{\upsilon}
\newcommand{\e}{\epsilon}
\newcommand{\tac}{\textasteriskcentered}

%mahesh macros
\newcommand{\tm}{\textrm}

\newcommand{\bone}{\mathbf{1}}
\newcommand{\sign}{\mathrm{sign}}
\newcommand{\eps}{\varepsilon}
\newcommand{\textui}[1]{\underline{\textit{#1}}}

%\newcommand{\ov}{\overline}
%\newcommand{\un}{\underline}
\newcommand{\fin}{\mathrm{fin}}
\newcommand{\nl}{\newline \mbox{}}
\newcommand{\h}[1]{\hspace{#1pt}}
\newcommand{\mtext}[1]{\hspace{6pt}\text{#1}\hspace{6pt}}
\newcommand{\lnorm}{\left\lVert}
\newcommand{\rnorm}{\right\rVert}
\newcommand{\ds}{\displaystyle}
\newcommand{\ts}{\textstyle}
\newcommand{\sfrac}[2]{{}^{#1}\mskip -5mu/\mskip -3mu_{#2}}

% ============================================================
% Category-style operators and contoured variants
% ============================================================

\DeclareMathOperator{\Sets}{S \mkern1.04mu e \mkern1.04mu t \mkern1.04mu s}
    \newcommand{\cSets}{\scalebox{1.02}{\contour{black}{$\Sets$}}}
    
\DeclareMathOperator{\Groups}{G \mkern1.04mu r \mkern1.04mu o \mkern1.04mu u \mkern1.04mu p \mkern1.04mu s}
    \newcommand{\cGroups}{\scalebox{1.02}{\contour{black}{$\Groups$}}}

\DeclareMathOperator{\TTop}{T \mkern1.04mu o \mkern1.04mu p}
    \newcommand{\cTop}{\scalebox{1.02}{\contour{black}{$\TTop$}}}

\DeclareMathOperator{\Htp}{H \mkern1.04mu t \mkern1.04mu p}
    \newcommand{\cHtp}{\scalebox{1.02}{\contour{black}{$\Htp$}}}

\DeclareMathOperator{\Mod}{M \mkern1.04mu o \mkern1.04mu d}
    \newcommand{\cMod}{\scalebox{1.02}{\contour{black}{$\Mod$}}}

\DeclareMathOperator{\Ab}{A \mkern1.04mu b}
    \newcommand{\cAb}{\scalebox{1.02}{\contour{black}{$\Ab$}}}

\DeclareMathOperator{\Rings}{R \mkern1.04mu i \mkern1.04mu n \mkern1.04mu g \mkern1.04mu s}
    \newcommand{\cRings}{\scalebox{1.02}{\contour{black}{$\Rings$}}}

\DeclareMathOperator{\ComRings}{C \mkern1.04mu o \mkern1.04mu m \mkern1.04mu R \mkern1.04mu i \mkern1.04mu n \mkern1.04mu g \mkern1.04mu s}
    \newcommand{\cComRings}{\scalebox{1.05}{\contour{black}{$\ComRings$}}}

\DeclareMathOperator{\hHom}{H \mkern1.04mu o \mkern1.04mu m}
    \newcommand{\cHom}{\scalebox{1.02}{\contour{black}{$\hHom$}}}

% ============================================================
% Mathcal
% ============================================================

\newcommand{\cA}{\mathcal{A}}
\newcommand{\cB}{\mathcal{B}}
\newcommand{\cC}{\mathcal{C}}
\newcommand{\cD}{\mathcal{D}}
\newcommand{\cE}{\mathcal{E}}
\newcommand{\cF}{\mathcal{F}}
\newcommand{\cG}{\mathcal{G}}
\newcommand{\cH}{\mathcal{H}}
\newcommand{\cI}{\mathcal{I}}
\newcommand{\cJ}{\mathcal{J}}
\newcommand{\cK}{\mathcal{K}}
\newcommand{\cL}{\mathcal{L}}
\newcommand{\cM}{\mathcal{M}}
\newcommand{\cN}{\mathcal{N}}
\newcommand{\cO}{\mathcal{O}}
\newcommand{\cP}{\mathcal{P}}
\newcommand{\cQ}{\mathcal{Q}}
\newcommand{\cR}{\mathcal{R}}
\newcommand{\cS}{\mathcal{S}}
\newcommand{\cT}{\mathcal{T}}
\newcommand{\cU}{\mathcal{U}}
\newcommand{\cV}{\mathcal{V}}
\newcommand{\cW}{\mathcal{W}}
\newcommand{\cX}{\mathcal{X}}
\newcommand{\cY}{\mathcal{Y}}
\newcommand{\cZ}{\mathcal{Z}}

% ============================================================
% Mathscrl% 
%============================================================

\newcommand{\scra}{\mathscr{a}}
\newcommand{\scrA}{\mathscr{A}}
\newcommand{\scrb}{\mathscr{b}}
\newcommand{\scrB}{\mathscr{B}}
\newcommand{\scrc}{\mathscr{c}}
\newcommand{\scrC}{\mathscr{C}}
\newcommand{\scrd}{\mathscr{d}}
\newcommand{\scrD}{\mathscr{D}}
\newcommand{\scre}{\mathscr{e}}
\newcommand{\scrE}{\mathscr{E}}
\newcommand{\scrf}{\mathscr{f}}
\newcommand{\scrF}{\mathscr{F}}
\newcommand{\scrg}{\mathscr{g}}
\newcommand{\scrG}{\mathscr{G}}
\newcommand{\scrh}{\mathscr{h}}
\newcommand{\scrH}{\mathscr{H}}
\newcommand{\scri}{\mathscr{i}}
\newcommand{\scrI}{\mathscr{I}}
\newcommand{\scrj}{\mathscr{j}}
\newcommand{\scrJ}{\mathscr{J}}
\newcommand{\scrk}{\mathscr{k}}
\newcommand{\scrK}{\mathscr{K}}
\newcommand{\scrl}{\mathscr{l}}
\newcommand{\scrL}{\mathscr{L}}
\newcommand{\scrm}{\mathscr{m}}
\newcommand{\scrM}{\mathscr{M}}
\newcommand{\scrn}{\mathscr{n}}
\newcommand{\scrN}{\mathscr{N}}
\newcommand{\scro}{\mathscr{o}}
\newcommand{\scrO}{\mathscr{O}}
\newcommand{\scrp}{\mathscr{p}}
\newcommand{\scrP}{\mathscr{P}}
\newcommand{\scrq}{\mathscr{q}}
\newcommand{\scrQ}{\mathscr{Q}}
\newcommand{\scrr}{\mathscr{r}}
\newcommand{\scrR}{\mathscr{R}}
\newcommand{\scrs}{\mathscr{s}}
\newcommand{\scrS}{\mathscr{S}}
\newcommand{\scrt}{\mathscr{t}}
\newcommand{\scrT}{\mathscr{T}}
\newcommand{\scru}{\mathscr{u}}
\newcommand{\scrU}{\mathscr{U}}
\newcommand{\scrv}{\mathscr{v}}
\newcommand{\scrV}{\mathscr{V}}
\newcommand{\scrw}{\mathscr{w}}
\newcommand{\scrW}{\mathscr{W}}
\newcommand{\scrx}{\mathscr{x}}
\newcommand{\scrX}{\mathscr{X}}
\newcommand{\scry}{\mathscr{y}}
\newcommand{\scrY}{\mathscr{Y}}
\newcommand{\scrz}{\mathscr{z}}
\newcommand{\scrZ}{\mathscr{Z}}

% ============================================================
% Mathfrak (missing \fi)
% ============================================================

\newcommand{\fa}{\mathfrak{a}}
\newcommand{\fA}{\mathfrak{A}}
\newcommand{\fb}{\mathfrak{b}}
\newcommand{\fB}{\mathfrak{B}}
\newcommand{\fc}{\mathfrak{c}}
\newcommand{\fC}{\mathfrak{C}}
\newcommand{\fd}{\mathfrak{d}}
\newcommand{\fD}{\mathfrak{D}}
\newcommand{\fe}{\mathfrak{e}}
\newcommand{\fE}{\mathfrak{E}}
\newcommand{\ff}{\mathfrak{f}}
\newcommand{\fF}{\mathfrak{F}}
\newcommand{\fg}{\mathfrak{g}}
\newcommand{\fG}{\mathfrak{G}}
\newcommand{\fh}{\mathfrak{h}}
\newcommand{\fH}{\mathfrak{H}}
\newcommand{\fI}{\mathfrak{I}}
\newcommand{\fj}{\mathfrak{j}}
\newcommand{\fJ}{\mathfrak{J}}
\newcommand{\fk}{\mathfrak{k}}
\newcommand{\fK}{\mathfrak{K}}
\newcommand{\fl}{\mathfrak{l}}
\newcommand{\fL}{\mathfrak{L}}
\newcommand{\fm}{\mathfrak{m}}
\newcommand{\fM}{\mathfrak{M}}
\newcommand{\fn}{\mathfrak{n}}
\newcommand{\fN}{\mathfrak{N}}
\newcommand{\fo}{\mathfrak{o}}
\newcommand{\fO}{\mathfrak{O}}
\newcommand{\fp}{\mathfrak{p}}
\newcommand{\fP}{\mathfrak{P}}
\newcommand{\fq}{\mathfrak{q}}
\newcommand{\fQ}{\mathfrak{Q}}
\newcommand{\fr}{\mathfrak{r}}
\newcommand{\fR}{\mathfrak{R}}
\newcommand{\fs}{\mathfrak{s}}
\newcommand{\fS}{\mathfrak{S}}
\newcommand{\ft}{\mathfrak{t}}
\newcommand{\fT}{\mathfrak{T}}
\newcommand{\fu}{\mathfrak{u}}
\newcommand{\fU}{\mathfrak{U}}
\newcommand{\fv}{\mathfrak{v}}
\newcommand{\fV}{\mathfrak{V}}
\newcommand{\fw}{\mathfrak{w}}
\newcommand{\fW}{\mathfrak{W}}
\newcommand{\fx}{\mathfrak{x}}
\newcommand{\fX}{\mathfrak{X}}
\newcommand{\fy}{\mathfrak{y}}
\newcommand{\fY}{\mathfrak{Y}}
\newcommand{\fz}{\mathfrak{z}}
\newcommand{\fZ}{\mathfrak{Z}}

% ============================================================
% Mathbf
% ============================================================

\newcommand{\bfA}{\mathbf{A}}
\newcommand{\bfB}{\mathbf{B}}
\newcommand{\bfC}{\mathbf{C}}
\newcommand{\bfD}{\mathbf{D}}
\newcommand{\bfE}{\mathbf{E}}
\newcommand{\bfF}{\mathbf{F}}
\newcommand{\bfG}{\mathbf{G}}
\newcommand{\bfH}{\mathbf{H}}
\newcommand{\bfI}{\mathbf{I}}
\newcommand{\bfJ}{\mathbf{J}}
\newcommand{\bfK}{\mathbf{K}}
\newcommand{\bfL}{\mathbf{L}}
\newcommand{\bfM}{\mathbf{M}}
\newcommand{\bfN}{\mathbf{N}}
\newcommand{\bfO}{\mathbf{O}}
\newcommand{\bfP}{\mathbf{P}}
\newcommand{\bfQ}{\mathbf{Q}}
\newcommand{\bfR}{\mathbf{R}}
\newcommand{\bfS}{\mathbf{S}}
\newcommand{\bfT}{\mathbf{T}}
\newcommand{\bfU}{\mathbf{U}}
\newcommand{\bfV}{\mathbf{V}}
\newcommand{\bfW}{\mathbf{W}}
\newcommand{\bfX}{\mathbf{X}}
\newcommand{\bfY}{\mathbf{Y}}
\newcommand{\bfZ}{\mathbf{Z}}

\newcommand{\bfa}{\mathbf{a}}
\newcommand{\bfb}{\mathbf{b}}
\newcommand{\bfc}{\mathbf{c}}
\newcommand{\bfd}{\mathbf{d}}
\newcommand{\bfe}{\mathbf{e}}
\newcommand{\bff}{\mathbf{f}}
\newcommand{\bfg}{\mathbf{g}}
\newcommand{\bfh}{\mathbf{h}}
\newcommand{\bfi}{\mathbf{i}}
\newcommand{\bfj}{\mathbf{j}}
\newcommand{\bfk}{\mathbf{k}}
\newcommand{\bfl}{\mathbf{l}}
\newcommand{\bfm}{\mathbf{m}}
\newcommand{\bfn}{\mathbf{n}}
\newcommand{\bfo}{\mathbf{o}}
\newcommand{\bfp}{\mathbf{p}}
\newcommand{\bfq}{\mathbf{q}}
\newcommand{\bfr}{\mathbf{r}}
\newcommand{\bfs}{\mathbf{s}}
\newcommand{\bft}{\mathbf{t}}
\newcommand{\bfu}{\mathbf{u}}
\newcommand{\bfv}{\mathbf{v}}
\newcommand{\bfw}{\mathbf{w}}
\newcommand{\bfx}{\mathbf{x}}
\newcommand{\bfy}{\mathbf{y}}
\newcommand{\bfz}{\mathbf{z}}

% ============================================================
% Blackboard bold
% ============================================================

\newcommand{\bbA}{\mathbb{A}}
\newcommand{\bbB}{\mathbb{B}}
\newcommand{\bbC}{\mathbb{C}}
\newcommand{\bbD}{\mathbb{D}}
\newcommand{\bbE}{\mathbb{E}}
\newcommand{\bbF}{\mathbb{F}}
\newcommand{\bbG}{\mathbb{G}}
\newcommand{\bbH}{\mathbb{H}}
\newcommand{\bbI}{\mathbb{I}}
\newcommand{\bbJ}{\mathbb{J}}
\newcommand{\bbK}{\mathbb{K}}
\newcommand{\bbL}{\mathbb{L}}
\newcommand{\bbM}{\mathbb{M}}
\newcommand{\bbN}{\mathbb{N}}
\newcommand{\bbO}{\mathbb{O}}
\newcommand{\bbP}{\mathbb{P}}
\newcommand{\bbQ}{\mathbb{Q}}
\newcommand{\bbR}{\mathbb{R}}
\newcommand{\bbS}{\mathbb{S}}
\newcommand{\bbT}{\mathbb{T}}
\newcommand{\bbU}{\mathbb{U}}
\newcommand{\bbV}{\mathbb{V}}
\newcommand{\bbW}{\mathbb{W}}
\newcommand{\bbX}{\mathbb{X}}
\newcommand{\bbY}{\mathbb{Y}}
\newcommand{\bbZ}{\mathbb{Z}}

\newcommand{\Bmu}{\mbox{$\raisebox{-0.59ex}
  {$l$}\hspace{-0.18em}\mu\hspace{-0.88em}\raisebox{-0.98ex}{\scalebox{2}
  {$\color{white}.$}}\hspace{-0.416em}\raisebox{+0.88ex}
  {$\color{white}.$}\hspace{0.46em}$}{}}  %need graphicx and xcolor. this produces blackboard bold mu 

% ============================================================
% Greek-letter aliases
% ============================================================

\newcommand{\jota}{\jmath}

% ============================================================
% Matrices
% ============================================================

\newcommand{\bmat}{\left( \begin{matrix}}
\newcommand{\emat}{\end{matrix} \right)}

\newcommand{\bbmat}{\left[ \begin{matrix}}
\newcommand{\ebmat}{\end{matrix} \right]}

\newcommand{\pmat}{\left( \begin{smallmatrix}}
\newcommand{\epmat}{\end{smallmatrix} \right)}

\newcommand{\mat}[4]{\begin{pmatrix}{#1}&{#2}\\{#3}&{#4}\end{pmatrix}}

% ============================================================
% Residues and averaged integrals
% ============================================================

\newcommand{\res}[1]{\underset{#1}{\RES}\,}

\makeatletter
\newcommand*{\mint}[1]{%
  % #1: overlay symbol
  \mint@l{#1}{}%
}
\newcommand*{\mint@l}[2]{%
  % #1: overlay symbol
  % #2: limits
  \@ifnextchar\limits{%
    \mint@l{#1}%
  }{%
    \@ifnextchar\nolimits{%
      \mint@l{#1}%
    }{%
      \@ifnextchar\displaylimits{%
        \mint@l{#1}%
      }{%
        \mint@s{#2}{#1}%
      }%
    }%
  }%
}
\newcommand*{\mint@s}[2]{%
  % #1: limits
  % #2: overlay symbol
  \@ifnextchar_{%
    \mint@sub{#1}{#2}%
  }{%
    \@ifnextchar^{%
      \mint@sup{#1}{#2}%
    }{%
      \mint@{#1}{#2}{}{}%
    }%
  }%
}
\def\mint@sub#1#2_#3{%
  \@ifnextchar^{%
    \mint@sub@sup{#1}{#2}{#3}%
  }{%
    \mint@{#1}{#2}{#3}{}%
  }%
}
\def\mint@sup#1#2^#3{%
  \@ifnextchar_{%
    \mint@sup@sub{#1}{#2}{#3}%
  }{%
    \mint@{#1}{#2}{}{#3}%
  }%
}
\def\mint@sub@sup#1#2#3^#4{%
  \mint@{#1}{#2}{#3}{#4}%
}
\def\mint@sup@sub#1#2#3_#4{%
  \mint@{#1}{#2}{#4}{#3}%
}
\newcommand*{\mint@}[4]{%
  % #1: \limits, \nolimits, \displaylimits
  % #2: overlay symbol: -, =, ...
  % #3: subscript
  % #4: superscript
  \mathop{}%
  \mkern-\thinmuskip
  \mathchoice{%
    \mint@@{#1}{#2}{#3}{#4}%
        \displaystyle\textstyle\scriptstyle
  }{%
    \mint@@{#1}{#2}{#3}{#4}%
        \textstyle\scriptstyle\scriptstyle
  }{%
    \mint@@{#1}{#2}{#3}{#4}%
        \scriptstyle\scriptscriptstyle\scriptscriptstyle
  }{%
    \mint@@{#1}{#2}{#3}{#4}%
        \scriptscriptstyle\scriptscriptstyle\scriptscriptstyle
  }%
  \mkern-\thinmuskip
  \int#1%
  \ifx\\#3\\\else_{#3}\fi
  \ifx\\#4\\\else^{#4}\fi  
}
\newcommand*{\mint@@}[7]{%
  % #1: limits
  % #2: overlay symbol
  % #3: subscript
  % #4: superscript
  % #5: math style
  % #6: math style for overlay symbol
  % #7: math style for subscript/superscript
  \begingroup
    \sbox0{$#5\int\m@th$}%
    \sbox2{$#5\int_{}\m@th$}%
    \dimen2=\wd0 %
    % => \dimen2 = width of \int
    \let\mint@limits=#1\relax
    \ifx\mint@limits\relax
      \sbox4{$#5\int_{\kern1sp}^{\kern1sp}\m@th$}%
      \ifdim\wd4>\wd2 %
        \let\mint@limits=\nolimits
      \else
        \let\mint@limits=\limits
      \fi
    \fi
    \ifx\mint@limits\displaylimits
      \ifx#5\displaystyle
        \let\mint@limits=\limits
      \fi
    \fi
    \ifx\mint@limits\limits
      \sbox0{$#7#3\m@th$}%
      \sbox2{$#7#4\m@th$}%
      \ifdim\wd0>\dimen2 %
        \dimen2=\wd0 %
      \fi
      \ifdim\wd2>\dimen2 %
        \dimen2=\wd2 %
      \fi
    \fi
    \rlap{%
      $#5%
        \vcenter{%
          \hbox to\dimen2{%
            \hss
            $#6{#2}\m@th$%
            \hss
          }%
        }%
      $%
    }%
  \endgroup
}
\newcommand{\avg}{\mint{-}}
\makeatother

% ============================================================
% Comments and drafting helpers
% ============================================================

%Comments
\newcommand{\com}[1]{\vspace{5 mm}\par \noindent
\marginpar{\textsc{Comment}} \framebox{\begin{minipage}[c]{0.95
\textwidth} \tt #1 \end{minipage}}\vspace{5 mm}\par}

% ============================================================
% Arrows
% ============================================================

\newcommand{\hooktwoheadrightarrow}{%
  \hookrightarrow\mathrel{\mspace{-15mu}}\rightarrow
}

\makeatletter
\newcommand{\xhooktwoheadrightarrow}[2][]{%
  \lhook\joinrel
  \ext@arrow 0359\rightarrowfill@ {#1}{#2}%
  \mathrel{\mspace{-15mu}}\rightarrow
}
\makeatother

\makeatletter
\newcommand*{\da@rightarrow}{\mathchar"0\hexnumber@\symAMSa 4B }
\newcommand*{\da@leftarrow}{\mathchar"0\hexnumber@\symAMSa 4C }
\newcommand*{\xdashrightarrow}[2][]{%
  \mathrel{%
    \mathpalette{\da@xarrow{#1}{#2}{}\da@rightarrow{\,}{}}{}%
  }%
}
\newcommand{\xdashleftarrow}[2][]{%
  \mathrel{%
    \mathpalette{\da@xarrow{#1}{#2}\da@leftarrow{}{}{\,}}{}%
  }%
}
\newcommand*{\da@xarrow}[7]{%
  % #1: below
  % #2: above
  % #3: arrow left
  % #4: arrow right
  % #5: space left 
  % #6: space right
  % #7: math style 
  \sbox0{$\ifx#7\scriptstyle\scriptscriptstyle\else\scriptstyle\fi#5#1#6\m@th$}%
  \sbox2{$\ifx#7\scriptstyle\scriptscriptstyle\else\scriptstyle\fi#5#2#6\m@th$}%
  \sbox4{$#7\dabar@\m@th$}%
  \dimen@=\wd0 %
  \ifdim\wd2 >\dimen@
    \dimen@=\wd2 %   
  \fi
  \count@=2 %
  \def\da@bars{\dabar@\dabar@}%
  \@whiledim\count@\wd4<\dimen@\do{%
    \advance\count@\@ne
    \expandafter\def\expandafter\da@bars\expandafter{%
      \da@bars
      \dabar@ 
    }%
  }%  
  \mathrel{#3}%
  \mathrel{%   
    \mathop{\da@bars}\limits
    \ifx\\#1\\%
    \else
      _{\copy0}%
    \fi
    \ifx\\#2\\%
    \else
      ^{\copy2}%
    \fi
  }%   
  \mathrel{#4}%
}
\makeatother

% ============================================================
% Relation and delimiter spacing
% ============================================================

\renewcommand{\geq}{\geqslant}
\renewcommand{\leq}{\leqslant}
\newcommand{\midd}{\hspace{4pt}\middle|\hspace{4pt}}

% ============================================================
% Left Right Updated Deliminators
% ============================================================

\makeatletter
\NewDocumentCommand\Left{O{}}{%
    \IfValueTF{#1}{\notblank{#1}{\mathopen\bBigg@{#1}}{\left}}{\left}}
\NewDocumentCommand\Right{O{}}{%
    \IfValueTF{#1}{\notblank{#1}{\mathopen\bBigg@{#1}}{\right}}{\right}}
\makeatother

% ============================================================
% Restrictions
% ============================================================

\newcommand{\littletaller}{\mathchoice{\vphantom{\big|}}{}{}{}}
\newcommand\restr[2]{{% we make the whole thing an ordinary symbol
  \left.\kern-\nulldelimiterspace % automatically resize the bar with \right
  #1 % the function
  \littletaller % pretend it's a little taller at normal size
  \right|_{#2} % this is the delimiter
  }}

% ============================================================
% Chapter title formatting
% ============================================================

\newcommand{\chnum}{\titleformat
{\chapter} % command
[display] % shape
{\centering} % format
{\Huge \color{black} \shadowbox{\thechapter}} % label
{-0.5em} % sep (space between the number and title)
{\LARGE \color{black} \underline} % before-code
}

\newcommand{\chunnum}{\titleformat
{\chapter} % command
[display] % shape
{} % format
{} % label
{0em} % sep
{ \begin{flushright} \begin{tabular}{r}  \Huge \color{black}
} % before-code
[
\end{tabular} \end{flushright} \normalsize
] % after-code
}


\begin{document}
\homeworktitle
{Math 541} %class
{Week 2 Homework} %assignment
{Gianluca Crescenzo} %name

%%%%%%%%%%%%%%%%%%%%%%%%%%%%%%%%%%%
%%%%%%%%%%%%%%%%%%%%%%%%%%%%%%%%%%%
%%%%%%%%%%%%%%%%%%%%%%%%%%%%%%%%%%%
\begin{exercise}
	Let $a \in \bfR$. For $x \in \bfR^n$, define:
	\begin{align*}
		f(x) &= \begin{cases}
			|x|^{-a}, & 0 < |x| < 1 \\
			0, & \text{otherwise}
		\end{cases} \\
		g(x) &= \begin{cases}
			|x|^{-a}, & |x| > 1 \\
			0, & \text{otherwise}
		\end{cases}.
	\end{align*}
	For which $a$ is $f \in L^1(\bfR^n)$? For which $a$ is $f \in L^p(\bfR^n)$ for a given $1 \leq p < \infty$? Same questions for $g$.
\end{exercise}
	\begin{proof}
		We have:
		\begin{align*}
			\int_{\bfR^n}^{} |f(x)|^p dx
			& = \int_{B(0,1)}^{} |x|^{-ap}dx \\
			& = \int_{0}^{1} \int_{S^{n-1}}^{} r^{-ap}r^{n-1}d \omega dr \\
			& \leq C \int_{0}^{1} r^{n-1-ap}dr \\
			& = C \lim_{\delta \rightarrow 0} \left( \frac{1}{n-ap} - \frac{\delta^{n-ap}}{n-ap} \right).  
		\end{align*}
		This limit converges only when $n-ap > 0$. Thus $f \in L^p(\bfR^n)$ for $1 \leq p < \infty$ for all $a < \frac{n}{p}$. Now consider:
		\begin{align*}
			\int_{\bfR^n}^{} |g(x)|^p 
			& = \int_{\bfR^n \setminus B(0,1)}^{}|x|^{-ap} dx\\
			& = \int_{1}^{\infty} \int_{S^{n-1}}^{} r^{-ap}r^{n-1}d \omega dr \\
			& \leq C \int_{1}^{\infty} r^{n-1-ap}dr \\
			& = C \lim_{\delta \rightarrow \infty} \left( \frac{\delta^{n-ap}}{n-ap} - \frac{1}{n-ap}\right). 
		\end{align*}
		This limit converges only when $n-ap < 0$. Thus $g \in L^p(\bfR^n)$ for $1 \leq p < \infty$ for all $a > \frac{n}{p}$.
	\end{proof}
%%%%%%%%%%%%%%%%%%%%%%%%%%%%%%%%%%%
%%%%%%%%%%%%%%%%%%%%%%%%%%%%%%%%%%%
%%%%%%%%%%%%%%%%%%%%%%%%%%%%%%%%%%%
\newpage
\begin{exercise}
	Let $U$	be open in $\bfR^n$. Find a function $u \in L^1(U)$ which is unbounded on any open subset $V \subset U$. Hint: start with an enumeration of the points in $U$ with rational coordinates. Then make us of translates of the function $f$ (with an appropriate value for $a$) from the previous exercise.
\end{exercise}
	\begin{proof}
		Let $f$ be given in the previous exercise for any $0 < a < n$. Let $\bfQ^n \cap U = \{q_k \mid k \in \bfN\}$ be an enumeration of the rationals. Define:
		\[
			u(x) := 
			\begin{cases}
				\sum_{k=1}^{\infty} 2^{-k}f(x-q_k), & x \not\in \bfQ^n \cap U \\
				0, & x \in \bfQ^n \cap U.
			\end{cases}
		\] 
		Since $u(x) =\sum_{k=1}^{\infty} 2^{-k}f(x-q_k)$ a.e., it is enough to show that the series of functions alone is in $L^1(U)$. Note that:
		\begin{align*}
			\int_{U}^{} \sum_{k=1}^{\infty} 2^{-k}f(x-q_k) dx = \sum_{k=1}^{\infty} \int_{U}^{}f(x-q_k)dx. 
		\end{align*}
		By the translating the function $f$, we have:
		\begin{align*}
			\int_{U}^{}f(x-q_k)dx 
			& = \int_{U - q_k}^{}f(y)dy \\
			& \leq \int_{\bfR^n}^{} f(y)dy \\
			& = \lnorm f \rnorm_{L^1(\bfR^n)}.
		\end{align*}
		Hence:
		\begin{align*}
			\int_{U}^{} \sum_{k=1}^{\infty} 2^{-k}f(x-q_k) dx
			& \leq \sum_{k=1}^{\infty} 2^{-k}\lnorm f \rnorm_{L^1(\bfR^n)} \\
			& = \lnorm f \rnorm_{L^1(\bfR^n)} \sum_{k=1}^{\infty} 2^{-k} \\
			& < \infty.
		\end{align*}
		Thus $u \in L^1(U)$. Now let $V \subset U$ be any nonempty open set. Since $\bfQ^n$ is dense in $\bfR^n$, there exists some $q_m \in V$. Since $V$ is open, we can find points $x \in V \setminus\bfQ^n$ which are arbitrarily close to $q_m$. Hence:
		\begin{align*}
			\sum_{k = 1}^{\infty} 2^{-k}f(x-q_k) \geq 2^{-m}f(x-q_m).
		\end{align*}
		If $0 < |x-q_m| < 1$, then $f(x-q_m) = |x-q_m|^{-a}$. This finally gives:
		\[
			\sum_{k = 1}^{\infty} 2^{-k}f(x-q_k) \geq 2^{-m}|x-q_m|^{-a}.
		\] 
		As $x \rightarrow q_m$, the right hand side of the above inequality tends to infinity. Thus $u$ is unbounded on $V$.
	\end{proof}
%%%%%%%%%%%%%%%%%%%%%%%%%%%%%%%%%%%
%%%%%%%%%%%%%%%%%%%%%%%%%%%%%%%%%%%
%%%%%%%%%%%%%%%%%%%%%%%%%%%%%%%%%%%
\newpage
\begin{exercise}
	Show that the fundamental solution $\Phi$ to Laplace's equation has a finite integral over any bounded neighborhood of the origin.
\end{exercise}
	\begin{proof}
		Let $U \subset \bfR^n$ be a bounded neighborhood of the origin. Then there exists $r > 0$ such that $U \subseteq B(0,r)$. Hence:
		\[
			\int_{U}^{} \left| \Phi(x) \right| dx \leq \int_{B(0,r)}^{} \left| \Phi(x) \right| dx
		\]
		For $n=2$:
		\begin{align*}
			\int_{B(0,r)}^{} \left| \Phi(x) \right| dx
			&= c_1 \int_{B(0,r)}^{} |\log|x||dx \\
			&= c_1 \int_{0}^{r} \int_{0}^{2\pi} |r \log r|d \theta dr \\
			& \leq C \int_{0}^{r} |r \log r|dr \\
			&< \infty. \\
		\end{align*}
		For $n \geq 3$ :
		\begin{align*}
			\int_{B(0,r)}^{} \left| \Phi(x) \right| dx
			&= c_n \int_{B(0,r)}^{} |x|^{2-n}dx \\
			&= c_n \int_{0}^{r} \int_{S^{n-1}}^{} r^{2-n}r^{n-1}d \omega dr\\
			&\leq C \int_{0}^{r} r^{2-n}r^{n-1}dr \\
			&= C \int_{0}^{r} r dr\\
			&< \infty.  \\
		\end{align*}
		Hence by our first inequality, the fundamental solution to Laplace's equation has a finite integral over $U$.
	\end{proof}
%%%%%%%%%%%%%%%%%%%%%%%%%%%%%%%%%%%
%%%%%%%%%%%%%%%%%%%%%%%%%%%%%%%%%%%
%%%%%%%%%%%%%%%%%%%%%%%%%%%%%%%%%%%
\begin{exercise}
	For $i,j = 1,\ldots,n$, calculate $\Phi_{x_i}(x)$ and $\Phi_{x_i x_j}(x)$ for any $n\geq 2$. Use your answer to show:
	\[
		|D \Phi(x)| \leq \frac{C}{|x|^{n-1}} \text{ and } |D^2 \Phi(x)|\leq \frac{C}{|x|^n}.
	\] 
\end{exercise}
	\begin{proof}
		We will first show that $|D \Phi(x)| \leq \frac{C}{|x|^{n-1}}$. We proceed by cases on $n$. Case 1: $n = 2$. We have $\frac{\partial}{\partial x_i}\Phi(x) = -c_2 \frac{x_i}{|x|^2}$, hence $D \Phi(x) = -c_2 \frac{x}{|x|^2}$. Thus:
		\[
			|D \Phi(x)| \leq c_2 \frac{|x|}{|x|^2} \leq \frac{C}{|x|}.
		\] 
		Case 2: $n \geq 3$. We have $\frac{\partial}{\partial x_i}\Phi(x) = c_n(2-n) \frac{x_i}{|x|^n}$, hence $D \Phi(x) = c_n (2-n) \frac{x}{|x|^n}$. Hence:
		\[
			|D \Phi(x)| \leq c_n(2-n) \frac{|x|}{|x|^n} \leq \frac{C}{|x|^{n-1}}.
		\]

		Using the above work, we can now show that $|D^2 \Phi(x)| \leq \frac{C}{|x|^n}$. For $n=2$:
		\begin{align*}
			\frac{\partial}{\partial x_j}\Phi_{x_i}(x)
			& = -c_2 \frac{\partial}{\partial x_j} \frac{x_i}{|x|^2} \\
			& = -c_2 \frac{\partial}{\partial x_j} \frac{x_i}{|x|^2} \\
			&= \delta_{ij} |x|^{-2} - 2 x_i x_j |x|^{-4}  \\
		\end{align*}
		For $n \geq 3$:
		\begin{align*}
			\frac{\partial}{\partial x_j}\Phi_{x_i}(x)
			& = -c_n(2-n) \frac{\partial}{\partial x_j} \frac{x_i}{|x|^n} \\
			& = \delta_{ij}|x|^{-n} - nx_i x_j |x|^{-n-2}.
		\end{align*}
		In either case of dimension, we have:
		\begin{align*}
			\bigl|\delta_{ij} |x|^{-n}\bigr| &\leq \frac{1}{|x|^n}, \\
			\left| \frac{x_i x_j}{|x|^{n+2}} \right| &\leq \frac{|x|^2}{|x|^{n+2}} = \frac{1}{|x|^n}.
		\end{align*}
		Thus:
		\[
			|\Phi_{x_i x_j}(x)| \leq \frac{C}{|x|^n}.
		\]
		Since every entry of the Hessian matrix is bounded by $\frac{C}{|x|^n}$, we have for $n \geq 2$:
		\[
			|D^2 \Phi(x)| \leq \frac{C}{|x|^n}. \qedhere
		\] 
	\end{proof}
%%%%%%%%%%%%%%%%%%%%%%%%%%%%%%%%%%
%%%%%%%%%%%%%%%%%%%%%%%%%%%%%%%%%%
%%%%%%%%%%%%%%%%%%%%%%%%%%%%%%%%%%
\begin{exercise}
	Let $f \in C^2_c(\bfR^n)$ and $1 > \epsilon > 0$. Recall the following integrals that came up in our calculations in class:
	\begin{align*}
		I_\epsilon &= \int_{B(0,\epsilon)}^{} \Delta_y f(x-y)\Phi(y)dy\\
		J_\epsilon &= \int_{\partial B(0,\epsilon}^{} \frac{\partial f}{\partial \nu}(x-y) \Phi(y)dy.
	\end{align*}
	Show that for $n = 2$ we get:
	\begin{enumerate}[label = (\alph*),itemsep=1pt,topsep=3pt]
		\item $ |I_\epsilon| \leq C \epsilon^2 |\log \epsilon|$.
		\item $|J_\epsilon| \leq C \epsilon |\log \epsilon|$.
	\end{enumerate}
\end{exercise}
	\begin{proof}
		(a) We have:
		\begin{align*}
			|I_\epsilon|
			&= \left| \int_{B(0,\epsilon}^{} \Delta_yf(x-y) \Phi(y)dy \right|  \\
			&\leq \int_{B(0,\epsilon}^{} \left| \Delta_y f(x-y) \right| \Phi(y)dy   \\
		\end{align*}
		Since $f \in C^2_c(\bfR^2)$, we have that $\Delta_y f(x-y)$ is bounded. Hence:
		\begin{align*}
			|I_\epsilon|
			& \leq C \int_{B(0,\epsilon)}^{} \Phi(y)dy \\
			& = C \int_{B(0,\epsilon)}^{} -c_2 \log |y| dy \\
			& \leq C \int_{B(0,\epsilon)}^{} \log\left( \frac{1}{|y|} \right) dy \\ 
			&= C \int_{0}^{\epsilon} \int_{S^1}^{} r \log \left( \frac{1}{r} \right) d\omega dr \\
			&\leq C \int_{0}^{\epsilon} r \log \left( \frac{1}{r} \right) dr.
		\end{align*}
		This is an improper integral at $r = 0$. After making the necessary adjustments, integration by parts yields:
		\begin{align*}
			C \lim_{\delta \rightarrow 0} \int_{\delta}^{\epsilon}r \log \left( \frac{1}{r} \right) dr 
			& = C \lim_{\delta \rightarrow 0} \left( \frac{r^2}{2}\log \left( \frac{1}{r} \right) \Bigg|_\delta^\epsilon + \frac{1}{2}\int_{\delta}^{\epsilon} r^{}dr\right)  \\
			&= C\lim_{\delta \rightarrow 0} \left( \frac{r^2}{2}\log \left( \frac{1}{r} \right) + \frac{r^2}{4} \right)\Bigg|_\delta^\epsilon  \\
			&=C \lim_{\delta \rightarrow 0} \left( \frac{\epsilon^2}{2}\log\left( \frac{1}{\epsilon} \right) + \frac{\epsilon^2}{4} - \frac{\delta^2}{2}\log\left( \frac{1}{\delta} \right) - \frac{\delta^2}{4}\right) \\ 
			&= C \left(  \frac{\epsilon^2}{2}\log\left( \frac{1}{\epsilon} \right) + \frac{\epsilon^2}{4}\right) \\
			&= C \epsilon^2 \left( \frac{1}{2}|\log\epsilon| + \frac{1}{4} \right) \\
			&\leq C \epsilon^2 \left( \frac{1}{2}|\log\epsilon| + \frac{1}{4}|\log\epsilon| \right)\\
			& \leq C \epsilon^2 \left( \frac{1}{2} + \frac{1}{4} \right)|\log\epsilon| \\
			& \leq C \epsilon^2 |\log\epsilon|.
		\end{align*}

		(b) Observe that:
		\begin{align*}
			|J_\epsilon|
			&= \left| \int_{\partial B(0,\e)}^{} \frac{\partial f}{\partial \nu}(x-y) \Phi(y)dy  \right|  \\
			&= \int_{_\partial B(0,\epsilon) }^{} \left|\frac{\partial f}{\partial \nu}(x-y)\right| \Phi(y)dy \\
		\end{align*}
		Since $f \in C^2_c(\bfR)$, we have that $\frac{\partial f}{\partial \nu}(x-y)$ is bounded. Hence:
		\begin{align*}
			|J_\epsilon| \leq C \int_{\partial B(0,\epsilon)}^{} \Phi(y)dy 
			&= C \int_{\partial B(0,\epsilon)}^{} -c_2 \log |y|dy \\
			&\leq C \int_{\partial B(0,\epsilon)}^{} \log \left( \frac{1}{|y|} \right) dy \\
			&\leq C \int_{\partial B(0,\epsilon)}^{} |\log |y|| dy\\
			&\leq C |\log \epsilon| \int_{B(0,\epsilon)}^{} dy \\
			& = C |\log \epsilon| 2 \pi \epsilon \\
			& \leq C \epsilon |\log \epsilon|. \qedhere
		\end{align*}	
	\end{proof}
%%%%%%%%%%%%%%%%%%%%%%%%%%%%%%%%%%%
%%%%%%%%%%%%%%%%%%%%%%%%%%%%%%%%%%%
%%%%%%%%%%%%%%%%%%%%%%%%%%%%%%%%%%%
\begin{exercise}
	Recall the integral:
	\[
		L_\epsilon = \int_{\partial B(0,\epsilon)}^{} \left( D \Phi(y) \cdot \vec{\nu} \right) f(x-y)dS(y).
	\] 
	We showed that for $n \geq 3$ we get $L_\epsilon = \mint{-} f(x-y)dS(y)$. Show that we get the same thing for $n = 2$.
\end{exercise}
	\begin{proof}
		First consider:
		\begin{align*}
			\frac{\partial}{\partial y_i}\Phi(y)
			&= -c_2 \frac{\partial}{\partial y_i} \log |y| \\
			&= -c_2 \frac{y_i}{|y|^2}.
		\end{align*}
		Hence:
		\[
			D \Phi(y)=-c_2\frac{y}{|y|^2}.
		\] 
		But since we're integrating on $\partial B(0,\epsilon)$, this means $D \Phi(y) = -c_2 \frac{y}{\epsilon^2}$. It remains to compute $D\Phi(y)\cdot \vec{\nu}$. Since we are on the ball of radius epsilon, the vector $\vec{\nu}$ is a function of the points $y \in \partial B(0,\epsilon)$. In particular, $\vec{\nu} = \frac{y}{\epsilon}$. This gives: 
		\begin{align*}
			D \Phi(y) \cdot \vec{\nu}
			&= -c_2 \frac{y}{\epsilon^2} \cdot \frac{y}{\epsilon} \\
			&= -c_2 \frac{1}{\epsilon^3}(y \cdot y) \\
			&= -c_2 \frac{1}{\epsilon^3}|y|^2 \\
			&= -c_2 \epsilon^{-1}. \\
		\end{align*}
		Finally:
		\begin{align*}
			L_\epsilon
			& = \int_{\partial B(0,\epsilon)}^{} -c_2 \epsilon^{-1}f(x-y)dS(y) \\
			& = -c_2 \epsilon^{-1} \frac{2 \pi \epsilon}{2 \pi \epsilon} \int_{\partial B(0,\epsilon)}^{} f(x-y)dS(y) \\
			& = -c_2 \epsilon^{-1} 2 \pi \epsilon \mint{-}_{\partial B(0,\epsilon)}^{}f(x-y)dS(y) \\
			& = -\mint{-}_{\partial B(0,\epsilon)}^{}f(x-y)dS(y). \qedhere
		\end{align*}
	\end{proof}
%%%%%%%%%%%%%%%%%%%%%%%%%%%%%%%%%%%
%%%%%%%%%%%%%%%%%%%%%%%%%%%%%%%%%%%
%%%%%%%%%%%%%%%%%%%%%%%%%%%%%%%%%%%
\begin{exercise}
	Let $\epsilon > 0$.
	\begin{enumerate}[label = (\alph*),itemsep=1pt,topsep=3pt]
		\item For $x \in \bfR$, let:
			\[
				H(x) :=
				\begin{cases}
					0, & x < 0 \\
					1, & x \geq 0.
				\end{cases}
			\]
			Compute the function $\zeta := \eta_\epsilon \ast H$. We call $\zeta$ the cutoff function.

		\item Let $a < b$ be real numbers. Find a smooth function which is indentically 1 on $[a,b]$, and is identically 0 for $x \leq a - \epsilon$ and $x \geq b + \epsilon$.

		\item Now let $U,V \subset \bfR^n$ be open sets with $V \subset\subset U$. Show that there is a smooth function on $\bfR^n$ which is identically 1 on $V$, and identically 0 outisde $U$.
	\end{enumerate}
\end{exercise}
	\begin{proof}
		(a) We have:
\begin{align*}
\zeta(x) &= (\eta_\epsilon * H)(x) \\
&= \int_{\mathbb{R}} \eta_\epsilon(x-y)\,H(y)\,dy \\
&= \int_0^\infty \eta_\epsilon(x-y)\,dy \\
&= \int_{-\infty}^x \eta_\epsilon(t)\,dt \\[6pt]
&=
\begin{cases}
0, & x \le -\epsilon, \\[6pt]
\displaystyle \int_{-\epsilon}^{x} \frac{c}{\epsilon}\exp\!\left(\frac{1}{\left(\frac{t}{\epsilon}\right)^2 - 1}\right)\,dt,
& -\epsilon < x < \epsilon, \\[10pt]
1, & x \ge \epsilon.
\end{cases}
\end{align*}
	\end{proof}





%%%%%%%%%%%%%%%%%%%%%%%%%%%%%%%%%%%
%%%%%%%%%%%%%%%%%%%%%%%%%%%%%%%%%%%
%%%%%%%%%%%%%%%%%%%%%%%%%%%%%%%%%%%
\end{document}
